\documentclass[a4paper, 12pt]{article}

\usepackage[english, russian]{babel}
\usepackage[T1, T2A]{fontenc}
\usepackage[utf8]{inputenc}
\usepackage{geometry}
\usepackage{amsmath}
\usepackage{amsfonts}
\usepackage{booktabs}
\usepackage[dvipsnames]{xcolor}
\usepackage[colorlinks]{hyperref}
\usepackage{indentfirst}
\usepackage{listings}

\geometry{left=2cm, right=2cm, bottom=2cm, top=2cm}

\hypersetup{
    linkcolor=blue,
    filecolor=magenta,
    urlcolor=cyan
}


\begin{document}
\section*{Ответы на контрольные вопросы}

\subsection*{1. Линейное преобразование, векторное пространство, кольца/поля}

\textbf{Линейное преобразование} — отображение $L: V \to W$ между векторными пространствами, сохраняющее сложение и умножение на скаляр: $L(u+v) = L(u)+L(v)$ и $L(\alpha u) = \alpha L(u)$ для всех $u,v \in V$, $\alpha \in \mathbb{F}$.

\textbf{Векторное пространство} — множество $V$ с двумя операциями (сложение векторов и умножение на скаляр из поля $\mathbb{F}$), удовлетворяющее аксиомам замкнутости, ассоциативности, существования нулевого и обратного элементов, дистрибутивности.

\textbf{В задании 1 и 2} (шифр Хилла) работаем в кольце $\mathbb{Z}_{32}$ (целые числа по модулю 32). \textbf{В задании 3} (код Хэмминга) работаем в поле Галуа $\mathbb{F}_2$ (двоичное поле, где $1+1=0$).

\subsection*{2. Кодирование и декодирование методом Хилла}

\textbf{Кодирование:} текст разбивается на блоки-векторы $d$ (размер $k \times k$ матрицы-ключа $K$). Каждый блок умножается на $K$ в кольце $\mathbb{Z}_n$: $c = d \cdot K \pmod{n}$. Результат $c$ — зашифрованный блок.

\textbf{Декодирование:} $d = c \cdot K^{-1} \pmod{n}$, где $K^{-1}$ — обратная матрица. Обратная существует только при $\gcd(\det(K), n) = 1$.

\subsection*{3. Как находить обратную матрицу и почему определитель не должен иметь общих делителей с $n$}

\textbf{Алгоритм:} для квадратной матрицы $A$ в $\mathbb{Z}_n$:
\begin{enumerate}
    \item Вычисляем $\det(A) \pmod{n}$.
    \item Находим $\det(A)^{-1} \pmod{n}$ через расширенный алгоритм Евклида (если $\gcd(\det(A),n)=1$).
    \item Вычисляем матрицу алгебраических дополнений $C$.
    \item Транспонируем: $\text{adj}(A) = C^T$.
    \item $A^{-1} = \det(A)^{-1} \cdot \text{adj}(A) \pmod{n}$.
\end{enumerate}

\textbf{Почему $\gcd(\det(A),n)=1$:} если общий делитель $d>1$ существует, то $\det(A) \equiv 0 \pmod{d}$, и обратный элемент не существует (деление невозможно). Без обратной матрицы дешифрация невозможна.

\subsection*{4. Rank матрицы и связь с определителем}

\textbf{Rank (ранг)} — максимальное число линейно независимых строк или столбцов матрицы. Для квадратной матрицы $A$ ($k \times k$):
\begin{itemize}
    \item $\det(A) \neq 0$ (в поле) $\Rightarrow \text{rank}(A) = k$ (полный ранг).
    \item $\det(A) = 0$ $\Rightarrow \text{rank}(A) < k$ — строки/столбцы линейно зависимы, обратной матрицы нет.
\end{itemize}
В кольце $\mathbb{Z}_n$ определитель может быть ненулевым, но не иметь обратного (если не взаимно прост с $n$); тогда ранг полный, но матрица необратима в этом кольце.

\subsection*{5. Алгоритм Евклида и применение в задании 1}

\textbf{Алгоритм Евклида} находит $\gcd(a,b)$ и коэффициенты $x,y$ такие, что $a \cdot x + b \cdot y = \gcd(a,b)$. \textbf{Расширенный алгоритм} позволяет найти обратный элемент $a^{-1} \pmod{n}$ при $\gcd(a,n)=1$.

\textbf{Применение в задании 1:} для каждой матрицы-ключа проверяем, что $\det(K)$ взаимно прост с $n=32$ (т.е. $\det(K)$ нечётно). Для дешифрации вычисляем $K^{-1} \pmod{32}$ с помощью расширенного алгоритма Евклида.

\subsection*{6. На какие части сообщения влияют ошибки при дешифрации? Зависимость?}

Одна ошибка в одном символе зашифрованного сообщения влияет на весь блок расшифрованного сообщения, так как дешифрация — умножение на обратную матрицу $K^{-1}$. Размер влияния пропорционален размеру блока: для ключа $2\times2$ — 2 символа, $3\times3$ — 3 символа, $4\times4$ — 4 символа.

\subsection*{7. Как подобрать матрицу ключа, чтобы одна ошибка соответствовала одной ошибке в дешифрованном сообщении?}

В матричном шифре Хилла это невозможно: матрица смешивает все символы блока. Любая ошибка в зашифрованном тексте распространяется на весь блок при дешифрации. Для ограничения влияния ошибки нужно использовать меньшие блоки или применять корректирующие коды (код Хэмминга в задании 3).

\subsection*{8. Линейно (не)зависимые векторы. Определитель матрицы с зависимыми столбцами}

\textbf{Линейно независимые векторы} $v_1,\dots,v_k$ — такие, что $\alpha_1 v_1 + \dots + \alpha_k v_k = 0$ влечёт $\alpha_i = 0$ для всех $i$. \textbf{Зависимые} — существуют ненулевые $\alpha_i$.

\textbf{Определитель} матрицы, составленной из линейно зависимых столбцов (или строк), равен 0. Это следует из мультилинейности и антисимметричности определителя.

\subsection*{9. Матричное представление линейной системы и решение}

Система $A x = b$ представляется матрицей коэффициентов $A$, вектором неизвестных $x$ и вектором правых частей $b$. \textbf{Решение:} если $A$ обратима ($\det(A) \neq 0$ в поле), то $x = A^{-1} b$. В кольце $\mathbb{Z}_n$ решение существует при $\gcd(\det(A),n)=1$; тогда $x = A^{-1} b \pmod{n}$, где $A^{-1}$ вычисляется через алгебраические дополнения.

\subsection*{10. Как составляются матрицы $G$ и $H$, за что они отвечают?}

\textbf{Матрица генератора $G$} ($4 \times 7$): $G = [I_4 \;|\; P]$, где $P$ ($4 \times 3$) — матрица проверки чётности. Отвечает за кодирование: $c = d \cdot G \pmod{2}$ (4 информационных бита $\to$ 7 кодовых бит).

\textbf{Проверочная матрица $H$} ($3 \times 7$): $H = [P^T \;|\; I_3]$. Отвечает за обнаружение и исправление ошибок: синдром $s = H \cdot c^T \pmod{2}$ указывает позицию ошибочного бита. Условие корректности: $G \cdot H^T \equiv 0 \pmod{2}$.

\subsection*{11. Span, Range (Im) и Nullspace (Ker). Совпадение пространств}

\textbf{Span} (линейная оболочка) векторов $v_1,\dots,v_k$ — множество всех их линейных комбинаций: $\{\alpha_1 v_1 + \dots + \alpha_k v_k\}$.

\textbf{Range (Im $A$)} — образ матрицы: $\{A x \mid x \in V\}$. \textbf{Nullspace (Ker $A$)} — ядро: $\{x \mid A x = 0\}$.

\textbf{Для $G$ ($4\times7$) в $\mathbb{F}_2$}: $\text{Im}(G)$ — 4-мерное подпространство кодовых слов в $\mathbb{F}_2^7$; $\text{Ker}(G) = \{0\}$ (строки независимы). Пространства не совпадают (разные размерности).

\textbf{Для $H$ ($3\times7$) в $\mathbb{F}_2$}: $\text{Im}(H) = \mathbb{F}_2^3$ (все возможные синдромы); $\text{Ker}(H)$ — кодовые слова без ошибок (4-мерное подпространство). Пространства совпадают только для квадратных невырожденных матриц (например, $I_4$ или обратимые ключи из задания 1), где $\text{Ker}(A) = \{0\}$ и $\text{Im}(A) = \mathbb{F}_2^k$ — полное пространство.

\subsection*{12. Сколько ошибок исправляет код Хэмминга (7,4)? Как найти, устранить и восстановить?}

Код Хэмминга $(7,4)$ исправляет \textbf{1 ошибку в каждом 7-битном блоке}. В сообщении из 5 блоков (35 бит) при распределении ошибок по разным блокам можно исправить до 5 ошибок.

\textbf{Нахождение ошибки:} вычисляем синдром $s = H \cdot c^T \pmod{2}$. Синдром $s \in \mathbb{F}_2^3$ кодирует позицию ошибочного бита (например, $s = (101)_2 \to$ бит 2).

\textbf{Устранение:} инвертируем бит в позиции, указанной синдромом.

\textbf{Восстановление сообщения:} из исправленного 7-битного блока извлекаем первые 4 информационных бита. Собираем 20 бит в 4 буквы по таблице алфавита (32 буквы, 5 бит на букву).


\end{document}
